%==============================================================================
%  Standalone note: Loss of coercivity-norm control for the OSGS method
%  in the reaction-dominated (large Da) regime.
%
%  Notation is kept compatible with the manuscript
%  "A stabilized finite element method for incompressible, inertial flows
%   in inhomogeneous porous media" (Casas, Gonzalez-Usua, Codina, de-Pouplana).
%
%  The macros below replicate the relevant ones defined inline in that paper
%  (\ViscProj, \Reyn, \Damk, \triplenorm, \dualpairing, \innerprod, ...), so the
%  equations can be lifted into the manuscript with no changes.  If you compile
%  this inside the paper itself, simply delete the PREAMBLE-DUPLICATE block.
%==============================================================================
\documentclass[11pt]{article}

\usepackage[utf8]{inputenc}
\usepackage[margin=1in]{geometry}
\usepackage{amsmath,amssymb,amsthm,mathtools}
\usepackage{bm}
\usepackage{xcolor}
\usepackage[colorlinks=true,allcolors=blue]{hyperref}
\usepackage{cleveref}

%================= PREAMBLE-DUPLICATE (mirrors the manuscript) =================
% Dimensionless numbers
\newcommand{\Reyn}{\operatorname{\mathit{Re}}}
\newcommand{\Damk}{\operatorname{\mathit{Da}}}

% Projection operators used in the viscous term (deviatoric/symmetric split)
\newcommand{\DPi}{\overset{\scriptscriptstyle{\text{D}}}{\Pi}}
\newcommand{\SPi}{\overset{\scriptscriptstyle{\text{S}}}{\Pi}}
\newcommand{\ViscProj}{\overset{\scriptscriptstyle{\text{DS}}}{\Pi}}

% Triple norm
\newcommand{\triplenorm}[1]{{\left\vert\kern-0.25ex\left\vert\kern-0.25ex\left\vert #1
\right\vert\kern-0.25ex\right\vert\kern-0.25ex\right\vert}}

% Inner-product / duality pairings (simplified, comma-separated arguments)
\newcommand{\innerprod}[1]{\left(#1\right)}
\newcommand{\dualpairing}[1]{\left\langle #1\right\rangle}

% Theorem environments
\theoremstyle{plain}
\newtheorem{lemma}{Lemma}
\newtheorem{proposition}{Proposition}
\theoremstyle{remark}
\newtheorem{remark}{Remark}
%============================ end PREAMBLE-DUPLICATE ===========================

% Shorthands local to this note
\newcommand{\uh}{\boldsymbol{u}_h}
\newcommand{\ph}{p_h}
\newcommand{\Uh}{U_h}
\newcommand{\Vh}{V_h}
\newcommand{\Wh}{W_h}
\newcommand{\Th}{\mathcal{T}_h}
\newcommand{\Xh}{\mathcal{X}_{h0}}
\newcommand{\Pio}{\Pi_{\boldsymbol{\tau}h}}            % FE projection in the tau-inner product
\newcommand{\Porth}{\widetilde{\Pi}}                  % subscale-space projection
\newcommand{\siga}{\widetilde{\sigma}_\alpha}
\newcommand{\tauNS}{\tau_{1,\text{NS}}}
\newcommand{\normh}[1]{\left\Vert #1\right\Vert_h}
\newcommand{\norm}[1]{\left\Vert #1\right\Vert}

\title{Why the orthogonal subscales lose coercivity-norm control \\[2pt]
in the reaction-dominated limit}
\author{}
\date{}

\begin{document}
\maketitle

\noindent
\textit{Scope.}\; This note isolates, in a setting and notation compatible with the
manuscript, the mechanism by which the coercivity proof of the ASGS method fails
to carry over to the OSGS method as the Damk\"ohler number grows. The conclusion
is that the failure is structural for a constant reaction coefficient---it is not
an artifact of the proof technique---but that it degrades the \emph{provable
coercivity norm}, not the \emph{optimal convergence rate}, which both methods
retain. We close by recording the precedents for this observation in the
literature and the sense in which the present framing is new.

%------------------------------------------------------------------------------
\section{Setting and the common stabilized form}
%------------------------------------------------------------------------------
We work in the linearized, constant-coefficient setting of the stability analysis
of the manuscript: uniform viscosity $\nu$, a uniform isotropic resistance
$\boldsymbol{\sigma}=\sigma\boldsymbol{1}$, a given advection field $\boldsymbol{a}$
with $\nabla\!\cdot(\alpha\boldsymbol{a})=0$, and a porosity field well resolved by
the mesh. Writing $\Uh=[\uh;\ph]\in\Xh$ for the finite element unknowns, both
methods of interest share the single stabilized bilinear form
%
\begin{equation}\label{eq:BS}
  B_\text{S}(\boldsymbol{a},\Uh,\Vh)
  = B(\boldsymbol{a},\Uh,\Vh)
  - \sum_{K\in\Th}\dualpairing{\mathcal{L}^*\Vh,\ \boldsymbol{\tau}_K\,\Porth\!\left[\mathcal{L}\Uh\right]}_K ,
\end{equation}
%
differing \emph{only} in the choice of the projection $\Porth$ onto the subscale
space $\widetilde{\mathcal{X}}$~\cite{codina2018variational,hughes2007variational}:
%
\begin{equation}\label{eq:choices}
  \text{ASGS:}\quad \Porth=\mathcal{I},
  \qquad\qquad
  \text{OSGS:}\quad \Porth=\mathcal{I}-\Pio ,
\end{equation}
%
where $\Pio$ is the projection onto $\Xh$ associated with the (possibly
$\boldsymbol{\tau}$-weighted) inner product, the simplest choice being the
$L^2$-projection~\cite{codina2002stabilized}. Here $\mathcal{L}$ is the
momentum--continuity operator of the linearized problem and $\mathcal{L}^*$ its
formal adjoint; in the reaction-dominated regime the only piece of either operator
that survives is the zeroth-order (reactive) one, so that
%
\begin{equation}\label{eq:reactive-part}
  \mathcal{L}\Uh\big|_{\text{react}} = \sigma\uh,
  \qquad
  \mathcal{L}^*\Vh\big|_{\text{react}} = \sigma\boldsymbol{v}_h .
\end{equation}
%
Recall the stabilization parameter and the reactive control quantity of the
manuscript,
%
\begin{equation}\label{eq:tau1-siga}
  \tau_1=\frac{1}{\alpha_K\tauNS^{-1}+\sigma},
  \qquad
  \siga\coloneqq\frac{\tauNS^{-1}\,\sigma}{\tauNS^{-1}+\sigma/\alpha_K},
  \qquad
  \tauNS=\Big(c_1\tfrac{\nu}{h^2}+c_2\tfrac{|\boldsymbol{w}|}{h}\Big)^{-1},
\end{equation}
%
with the reaction-dominated asymptotics ($\Damk_h\to\infty$, $h\to0$)
%
\begin{equation}\label{eq:asymp}
  \tau_1\sim\frac{\alpha_K}{\sigma},
  \qquad
  \siga\sim\alpha_K(1+\Reyn_h)\frac{\nu}{h^2}
  \;\sim\;\frac{\alpha_K(1+\Reyn_h)}{\Damk_h}\,\frac{\sigma}{1},
\end{equation}
%
where the last form makes explicit that $\siga/\sigma\sim(1+\Reyn_h)/\Damk_h\to0$.

%------------------------------------------------------------------------------
\section{The ASGS coercivity argument, reduced to its reactive content}
%------------------------------------------------------------------------------
Testing \cref{eq:BS} with $\Vh=\Uh$ and $\Porth=\mathcal{I}$ reproduces the
manuscript's identity
%
\begin{align}\label{eq:ASGS-expand}
  B_\text{S}(\boldsymbol{a},\Uh,\Uh)
  &= 2\nu\big\|\alpha^{1/2}\ViscProj\nabla\uh\big\|^2
   + \big\|\sigma^{1/2}\uh\big\|^2
   + \varepsilon\big\|\ph\big\|^2 \notag\\
  &\quad
   + \big\|\tau_1^{1/2}\alpha X(\Uh)\big\|_h^2
   - \big\|\tau_1^{1/2}\big(2\nabla\!\cdot(\alpha\nu\ViscProj\nabla\uh)-\sigma\uh\big)\big\|_h^2
   + \big\|\tau_2^{1/2}\nabla\!\cdot(\alpha\uh)\big\|_h^2
   - \varepsilon^2\big\|\tau_2^{1/2}\ph\big\|_h^2 ,
\end{align}
%
with $X(\Uh)\coloneqq\boldsymbol{a}\cdot\nabla\uh+\nabla\ph$. The decisive feature
is purely algebraic. Expanding the squared norm of the difference on the second
line produces a \emph{cross term} and, crucially, a \emph{positive} reactive
square:
%
\begin{equation}\label{eq:positive-square}
  -\big\|\tau_1^{1/2}\big(2\nabla\!\cdot(\alpha\nu\ViscProj\nabla\uh)-\sigma\uh\big)\big\|_h^2
  = -\big\|\tau_1^{1/2}\,2\nabla\!\cdot(\alpha\nu\ViscProj\nabla\uh)\big\|_h^2
    \;\underbrace{-\;\big\|\tau_1^{1/2}\sigma\uh\big\|_h^2}_{\text{from }\mathcal{L}^*\text{ paired with }\mathcal{L}}
    \;+\;2\dualpairing{2\tau_1\nabla\!\cdot(\nu\alpha\ViscProj\nabla\uh),\,\sigma\uh}_h .
\end{equation}
%
After the cross term is absorbed with Young's inequality and the inverse estimate
$\|\nabla\!\cdot(\alpha\ViscProj\nabla\uh)\|_K\le C_\text{inv}h^{-1}
\|\alpha^{1/2}\ViscProj\nabla\uh\|_K$~\cite{brenner2008mathematical}, the surviving
coefficient of $\|\uh\|_K^2$ is controlled from below by $C\,\siga$, giving the
manuscript's bound
%
\begin{equation}\label{eq:ASGS-lemma}
  B_\text{S}(\boldsymbol{a},\Uh,\Uh)\ \ge\ C\,\triplenorm{\Uh}^2,
  \qquad
  \triplenorm{\Uh}^2
  = \nu\big\|\alpha^{1/2}\ViscProj\nabla\uh\big\|^2
  + \big\|\siga^{1/2}\uh\big\|^2
  + \varepsilon\big\|\ph\big\|^2
  + \big\|\tau_1^{1/2}\alpha X(\Uh)\big\|_h^2
  + \big\|\tau_2^{1/2}\nabla\!\cdot(\alpha\uh)\big\|_h^2 .
\end{equation}
%
The point to retain: the $\siga$-control of $\uh$ comes from the term
$-\|\tau_1^{1/2}\sigma\uh\|_h^2$ in \cref{eq:positive-square}, i.e.\ from the
stabilization \emph{pairing the reactive part of $\mathcal{L}^*\Vh$ with the
reactive part of $\mathcal{L}\Uh$}. This pairing is nonzero precisely because
$\Porth=\mathcal{I}$ leaves $\sigma\uh$ untouched.

%------------------------------------------------------------------------------
\section{The analogous OSGS argument and where it breaks}
%------------------------------------------------------------------------------
Repeating the computation with $\Porth=\mathcal{I}-\Pio$ replaces the stabilizing
term in \cref{eq:BS}, when tested with $\Uh$, by the seminorm of the
\emph{fine-scale} residual only,
%
\begin{equation}\label{eq:OSGS-stab}
  \sum_K\dualpairing{\mathcal{L}^*\Uh,\ \boldsymbol{\tau}_K(\mathcal{I}-\Pio)\mathcal{L}\Uh}_K
  = \sum_K\big\|\boldsymbol{\tau}_K^{1/2}(\mathcal{I}-\Pio)\mathcal{L}\Uh\big\|_K^2 \ \ge\ 0 ,
\end{equation}
%
which controls only the component of the residual orthogonal to $\Xh$. Now apply
\cref{eq:reactive-part}. Since $\uh\in\mathcal{V}_{h0}\subset\Xh$ and $\sigma$ is
constant, the reactive residual lies entirely in the finite element space, hence
%
\begin{equation}\label{eq:annihilation}
  \boxed{\;(\mathcal{I}-\Pio)\big[\sigma\uh\big]=\boldsymbol{0}\quad(\sigma\text{ constant}).\;}
\end{equation}
%
This is the same fact that, in the manuscript, justifies dropping the reactive
terms from the computed projection (their projection ``is exactly zero for
constant $\boldsymbol{\sigma}$''). Algorithmically convenient there, it is
analytically obstructive here: the positive reactive square that ASGS extracted in
\cref{eq:positive-square} is, for OSGS, \emph{identically zero in the
reaction-dominated limit}. The only remaining source of $L^2$-control on $\uh$ is
the bare Galerkin reactive term $\|\sigma^{1/2}\uh\|^2$ already present in
\cref{eq:ASGS-expand}, which carries the coefficient $\sigma$ rather than the
stabilization-enhanced $\siga$. Collecting this:

\begin{proposition}[Reactive coercivity gap]\label{prop:gap}
Under the constant-coefficient assumptions above, the largest coefficient
$c_{\uh}$ for which an estimate of the form
$B_\text{S}(\boldsymbol{a},\Uh,\Uh)\ge c_{\uh}\,\|\uh\|^2+(\text{other nonnegative terms})$
is provable from \cref{eq:BS} satisfies, as $\Damk_h\to\infty$,
\begin{equation}\label{eq:gap}
  c_{\uh}^{\,\mathrm{ASGS}}\sim C\,\siga\sim\alpha_K(1+\Reyn_h)\frac{\nu}{h^2},
  \qquad
  c_{\uh}^{\,\mathrm{OSGS}}\sim \sigma,
  \qquad\Longrightarrow\qquad
  \frac{c_{\uh}^{\,\mathrm{OSGS}}}{c_{\uh}^{\,\mathrm{ASGS}}}\sim\frac{\Damk_h}{1+\Reyn_h}\to\infty,
\end{equation}
i.e.\ the OSGS coercivity argument controls $\uh$ in a norm weaker than
$\triplenorm{\cdot}$ by a relative factor $\sim(1+\Reyn_h)/\Damk_h$ on the
$\siga$-block. No reorganization of the proof recovers the $\siga$-control from
\cref{eq:OSGS-stab} alone, because the relevant residual component is zero by
construction \eqref{eq:annihilation}.
\end{proposition}

\begin{remark}[Why only the reactive term is affected]
The annihilation \eqref{eq:annihilation} is special to the reaction operator,
which acts as a (scaled) identity and therefore cannot map $\uh$ out of $\Xh$. By
contrast, $\boldsymbol{a}\cdot\nabla\uh$ and $\nabla\ph$ generically possess
nonzero components orthogonal to $\Xh$, so the convective and pressure-gradient
stabilizations---the $\tau_1\|\alpha X(\Uh)\|_h^2$ and $\tau_2$-divergence blocks
of $\triplenorm{\cdot}$---survive the orthogonal projection intact. The deficiency
is thus confined to the zeroth-order term, exactly the term that
LBB/convection-oriented OSGS analyses never needed to control.
\end{remark}

%------------------------------------------------------------------------------
\section{Why optimal convergence rates are nevertheless retained}
%------------------------------------------------------------------------------
The gap of \cref{prop:gap} is a statement about the \emph{provable coercivity
norm}, and must not be read as a loss of accuracy. Three observations, all
consistent with the manuscript's numerical evidence, explain why the convergence
\emph{rate} is unaffected.

\smallskip
\noindent\textbf{(i) Consistency is untouched.} The OSGS residual fed back through
$\Porth$ is still the true finite element residual; orthogonality changes only
which projection of it enters \cref{eq:BS}. The interpolation-error right-hand side
of the convergence theorem is therefore identical for both methods, and the
estimate
$
  \triplenorm{E_h}\lesssim\sum_K h_K^{-1}\big(\tau_{2,K}^{1/2}\mathrm{E}_{\text{int},K}(\boldsymbol{u})
  +\tau_{1,K}^{1/2}\mathrm{E}_{\text{int},K}(p)\big)
$
holds in the (smaller) OSGS norm with the same order.

\smallskip
\noindent\textbf{(ii) The Galerkin reaction term covers the velocity.} In the
reaction-dominated limit the velocity error is itself governed by the shared term
$\|\sigma^{1/2}\boldsymbol{e}_u\|$. The very fact that kills the OSGS coercivity
contribution---$\sigma\uh$ being well represented in $\Xh$---also means this
Galerkin term represents the reaction accurately. This is the same mechanism the
manuscript invokes when it notes that the apparently weaker $\siga$-control ``turns
out not to be the case'' once the asymptotics are examined.

\smallskip
\noindent\textbf{(iii) Variable resistance escapes the cancellation.} For a
spatially varying $\sigma(\boldsymbol{x})$---as in the Darcy--Forchheimer closure
actually solved, where $\sigma$ depends on $\alpha(\boldsymbol{x})$ and on
$|\boldsymbol{u}|$---one has $\sigma\uh\notin\mathcal{V}_{h0}$ in general, so
$(\mathcal{I}-\Pio)[\sigma\uh]\neq\boldsymbol{0}$ and OSGS recovers a genuine,
nonzero reactive contribution. The degeneracy of \cref{prop:gap} is therefore the
\emph{worst case of an idealized constant-coefficient model problem}, not of the
formulation as implemented.

\smallskip
This reconciles the theory with the reported experiments, in which ASGS and OSGS
deliver essentially identical, optimal convergence rates across the
reaction-dominated columns of the parameter sweep.

%------------------------------------------------------------------------------
\section{Relation to the literature}
%------------------------------------------------------------------------------
The \emph{general principle} behind \cref{eq:annihilation}---that orthogonal
subscales provide no stabilization for residual components already representable in
the finite element space---is established in the variational multiscale
literature. It is stated explicitly by Gravenkamp and
Codina~\cite{gravenkamp2023dispersed}, who note that the stabilization vanishes
when the argument of the orthogonal projection lies in the finite element space,
and who trace exactly that mechanism for the \emph{convective} term. The
companion fact that the OSGS stabilized norm is constructed to control the
\emph{convective} term (with the reaction coefficient entering only through the
stabilization parameter, never as an independently controlled quantity) is visible
throughout the orthogonal-subscale analyses of the convection--diffusion--reaction
and Oseen problems~\cite{codina2002stabilized,codina-blasco2002cdr,codina2008analysis}.
Within the present manuscript, the constant-$\sigma$ vanishing is already used at
the algorithmic level (the reactive terms are excluded from the computed
projection because their projection is exactly zero).

What appears \emph{not} to have been stated in this form is the combination of:
(a) the explicit side-by-side coercivity comparison of \cref{eq:positive-square}
and \cref{eq:OSGS-stab}, identifying $-\|\tau_1^{1/2}\sigma\uh\|_h^2$ as precisely
the term ASGS keeps and OSGS cannot, with the gap quantified as $\sim1/\Damk_h$ in
\cref{eq:gap}; (b) the specialization to the reaction-dominated regime of the
\emph{porous} (generalized Navier--Stokes) problem, with the porosity-weighted
$\siga$; and (c) the reconciliation of items (i)--(iii) showing that the weaker
provable bound coexists with optimal observed rates, the variable-$\sigma$ remark
softening the result further. We therefore recommend presenting it not as a new
phenomenon but as making explicit, for the reaction-dominated regime, a known
limitation of orthogonal subscales.

%------------------------------------------------------------------------------
\section*{Suggested one-paragraph version for the manuscript}
%------------------------------------------------------------------------------
\noindent\textit{(Drop-in text, parallel to the existing remark on $\siga$.)}
\begin{quote}\small
It is worth noting that the coercivity argument above does not carry over verbatim
to the OSGS method in the reaction-dominated regime. The stabilizing term then
controls only the component of the residual orthogonal to $\Xh$, whereas for a
constant resistance $\sigma$ the reactive residual $\sigma\uh$ lies entirely in
$\mathcal{V}_{h0}$ and is annihilated by the orthogonal projection---the same
property that lets us exclude these terms from the computed projection. The
positive contribution $-\|\tau_1^{1/2}\sigma\uh\|_h^2$ that yields the
$\siga$-control in \cref{eq:ASGS-lemma} is therefore unavailable for OSGS, and the
velocity is controlled only through the bare Galerkin term $\|\sigma^{1/2}\uh\|^2$,
weaker by a factor $\sim(1+\Reyn_h)/\Damk_h$. This is a known feature of orthogonal
subscales~\cite{gravenkamp2023dispersed,codina2002stabilized}: no stabilization is
added for residual components already representable on the mesh. By consistency the
convergence \emph{rate} is unaffected---as our experiments confirm---and for the
spatially varying $\sigma(\boldsymbol{x})$ of the closure actually solved the
cancellation is no longer exact, so a nonzero reactive stabilization is in fact
recovered.
\end{quote}

%------------------------------------------------------------------------------
\begin{thebibliography}{99}
\itemsep2pt

\bibitem{codina2001stabilized}
R.~Codina,
\emph{A stabilized finite element method for generalized stationary
incompressible flows},
Comput. Methods Appl. Mech. Engrg. \textbf{190} (2001) 2681--2706.

\bibitem{codina2002stabilized}
R.~Codina,
\emph{Stabilized finite element approximation of transient incompressible flows
using orthogonal subscales},
Comput. Methods Appl. Mech. Engrg. \textbf{191} (2002) 4295--4321.

\bibitem{codina-blasco2002cdr}
R.~Codina, J.~Blasco,
\emph{Analysis of a stabilized finite element approximation of the transient
convection--diffusion--reaction equation using orthogonal subscales},
Comput. Vis. Sci. \textbf{4} (2002) 167--174.

\bibitem{codina2008analysis}
R.~Codina,
\emph{Analysis of a stabilized finite element approximation of the Oseen equations
using orthogonal subscales},
Appl. Numer. Math. \textbf{58} (2008) 264--283.

\bibitem{codina2018variational}
R.~Codina, S.~Badia, J.~Baiges, J.~Principe,
\emph{Variational multiscale methods in computational fluid dynamics},
in: Encyclopedia of Computational Mechanics, 2nd ed., Wiley, 2018, pp.~1--28.

\bibitem{hughes2007variational}
T.~J.~R. Hughes, G.~Sangalli,
\emph{Variational multiscale analysis: the fine-scale Green's function,
projection, optimization, localization, and stabilized methods},
SIAM J. Numer. Anal. \textbf{45} (2007) 539--557.

\bibitem{gravenkamp2023dispersed}
H.~Gravenkamp, R.~Codina, et al.,
\emph{A stabilized finite element method for modeling dispersed multiphase flows
using orthogonal subgrid scales},
2023. (Replace with the exact journal/volume/pages used in your bibliography.)

\bibitem{brenner2008mathematical}
S.~C. Brenner, L.~R. Scott,
\emph{The Mathematical Theory of Finite Element Methods},
3rd ed., Springer, 2008.

\end{thebibliography}

\end{document}
